% Created 2024-07-10 Wed 17:03
% Intended LaTeX compiler: pdflatex
\documentclass[11pt]{article}
\usepackage[utf8]{inputenc}
\usepackage[T1]{fontenc}
\usepackage{graphicx}
\usepackage{longtable}
\usepackage{wrapfig}
\usepackage{rotating}
\usepackage[normalem]{ulem}
\usepackage{amsmath}
\usepackage{amssymb}
\usepackage{capt-of}
\usepackage{hyperref}
\usepackage{chunkymacros}
\usepackage{mathtools}
\usepackage[a4paper,margin=2cm]{geometry}
\author{Euan Mendoza}
\date{\today}
\title{}
\hypersetup{
 pdfauthor={Euan Mendoza},
 pdftitle={},
 pdfkeywords={},
 pdfsubject={},
 pdfcreator={Emacs 29.4 (Org mode 9.8)}, 
 pdflang={English}}
\usepackage{biblatex}
\addbibresource{/home/emendoza/Documents/Research/references.bib}
\begin{document}

\tableofcontents

\section{Inbox}
\label{sec:org5a95c71}
\subsection{{\bfseries\sffamily TODO} Draft Report Version 1.}
\label{sec:org92e5479}
\subsubsection{{\bfseries\sffamily TODO} Introduction [0/4]}
\label{sec:orgb53d6fd}
\url{graph\_homomorphism.tex}

The \emph{Introduction} section should contain information about the computational complexity and motivation for writing a memetic algorithm. It should also give a brief introduction to the actual problem and a brief overview of the report structure.

\begin{itemize}
\item[{$\square$}] Introduce the problem.
\item[{$\square$}] Describe the problems complexity, why the need for a memetic algorithm?
\item[{$\square$}] Introduce some research questions around the problem.
\end{itemize}
\begin{enumerate}
\item {\bfseries\sffamily TODO} Preliminaries [0/4]
\label{sec:org9d14430}
\url{graph\_homomorphism.tex}

Setup the notation and introductory maths.

\begin{itemize}
\item[{$\square$}] Setup the standard notation.
\item[{$\square$}] Describe the groups and group actions.
\item[{$\square$}] Describe the graph notation.
\item[{$\square$}] Linear algebra notation.
\end{itemize}
\end{enumerate}
\subsubsection{{\bfseries\sffamily TODO} Literature Review/Problem and Background [0/5]}
\label{sec:org1e8b06c}
\url{graph\_homomorphism.tex}

We should probably give a more formal and complete definition to the problem here.
The literature review should give a general overview of prior related research and what the current state of the art approaches are.

\begin{itemize}
\item[{$\square$}] Formal statement of the problem, including some linear algebraic formulation as per discussions.
\item[{$\square$}] Background to Graph Isomorphism and Linear Algebraic Techniques. [0/2]
\begin{itemize}
\item[{$\square$}] Weisfeiler-Leman and practical graph isomorphism algorithms \autocite{WeisfeilerLeman1968} and see survey\autocite{GroheNeuen2021a}.
\item[{$\square$}] Linear Algebraic Analogues\autocite{LiQiao2017}\autocite{BrooksbankEtAl2019}\autocite{ChenEtAl2024}.
\end{itemize}
\item[{$\square$}] Background to the Graph Homomorphism/Subgraph Isomorphism Problem and current techniques. [0/2]
\begin{itemize}
\item[{$\square$}] Backtracking Techniques\autocite{JuttnerMadarasi2018,CarlettiEtAl2017}
\item[{$\square$}] Also wikipedia has some nice results listed. \href{https://en.wikipedia.org/wiki/Subgraph\_isomorphism\_problem}{Subgraph Isomorphism Problem Wikipedia Link}
\end{itemize}
\item[{$\square$}] Mention papers describing memetic approaches to the problem including \autocite{HaqueEtAl2019}
\end{itemize}
\subsection{{\bfseries\sffamily TODO} The Algorithm [0/2]}
\label{sec:org614c2c5}
\url{graph\_homomorphism.tex}

Here we describe the overall methodology used to construct the algorithm.

\begin{itemize}
\item[{$\square$}] Introduce the linear algebraic techniques used to construct an optimisation problem.
\item[{$\square$}] Discuss any considerations and enhancement techniques.
\end{itemize}
\subsection{{\bfseries\sffamily TODO} The Results [0/3]}
\label{sec:org1ce0691}
\url{graph\_homomorphism.tex}

Discuss and compare the results benchmarks to other implementations. What kind of graphs yield the best performance.

\begin{itemize}
\item[{$\square$}] Describe the optimal input landscape.
\item[{$\square$}] Compare with the other algorithm implementations.
\item[{$\square$}] What went bad or good and why?
\end{itemize}
\subsection{{\bfseries\sffamily TODO} The Conclusion}
\label{sec:org7f66cb9}
\url{graph\_homomorphism.tex}

Summary of overall paper.
\end{document}
